% !TEX root = ../main.tex
\chapter{Supplementary Material for the RPC Optimisation Study}

\label{chap:rpc_appendix}

This appendix provides additional technical details for the RPC
design optimisation study summarised in
Section~\ref{sec:rpc_optimisation}. It includes a detailed
description of the detector geometry, the multi-physics simulation
framework, and the prototype construction and performance
evaluation. The material presented here follows closely the results
published in Ref.~\cite{Mao2021JINST_RPC}.

\section{Detector Geometry and Spacer Configurations}

The studied glass RPC consists of two electrodes of size
$1\times1~\si{m^2}$, with thicknesses of
\SI{0.7}{\milli\metre} (anode) and
\SI{1.1}{\milli\metre} (cathode). The gas gap between the electrodes
is \SI{1.2}{\milli\metre} wide and is maintained by cylindrical
nylon spacers with a diameter of \SI{8}{\milli\metre}.

The chamber is enclosed by FR4 walls and sealed using silicone-based
glue. Gas is injected and extracted through pipes located at
opposite corners, with two inlet pipes and four outlet pipes to
avoid over-pressure effects. Gas distribution is further improved
using perforated \textsf{PEEK} gas guides positioned near the inlet
and outlet.

\begin{figure}[htbp]
    \centering
    \includegraphics[width=0.9\textwidth]{RPC_images/RPC.png}
    \caption{Schematic cross section of the glass RPC.}
    \label{fig:rpc_appendix_geometry}
\end{figure}

Two spacer configurations are considered: a reference design with
a regular $10\times10$ grid (100 spacers), and an optimised
shifted-spacer configuration with 76 spacers arranged in a staggered
pattern. The central region of the chamber is used for quantitative
comparison in order to minimise boundary effects.

\section{Multi-Physics Simulation Framework}

The optimisation study is based on coupled multi-physics simulations
performed using \textsc{COMSOL Multiphysics}. The model combines gas
flow, electrostatics and structural mechanics in a self-consistent
framework.

\subsection{Gas Flow Modelling}

The gas flow inside the chamber is described by the incompressible
Navier--Stokes equations:
\begin{equation}
\rho(\boldsymbol{u}\cdot\nabla)\boldsymbol{u}
= \nabla\cdot\left[-p\boldsymbol{I}
+ \mu(\nabla\boldsymbol{u}+(\nabla\boldsymbol{u})^{T})\right]
+ \boldsymbol{F},
\end{equation}
\begin{equation}
\nabla\cdot\boldsymbol{u} = 0.
\end{equation}

No-slip boundary conditions are applied at the chamber walls, and a
gas flow rate of \SI{1}{\liter\per\hour} is used, corresponding to
approximately 80\% gas renewal per hour.

\subsection{Electrostatic Modelling}

The electric field in the gas gap is described by:
\begin{equation}
\nabla\cdot\boldsymbol{D} = \rho_e, \qquad
\boldsymbol{E} = -\nabla V.
\end{equation}

\subsection{Structural Mechanics}

Mechanical deformation of the electrodes is computed using standard
elasticity equations. The forces originate from gas pressure and
electrostatic attraction between electrodes.

\subsection{Multi-Physics Coupling}

The coupling between gas flow, deformation and electrostatics is
implemented through pressure loading and electro-mechanical
interactions, ensuring a consistent evaluation of detector
behaviour under operating conditions.

\section{Detailed Simulation Results}

The performance of the two spacer configurations is evaluated by
comparing the gas-flow properties, mechanical deformation and
electric-field distributions obtained from the multi-physics
simulation. In each case, the reference configuration (left) and
the shifted-spacer configuration (right) are shown for direct
comparison.

\begin{figure}[htbp]
    \centering
    \includegraphics[width=0.48\textwidth]{RPC_images/Velocity1.png}
    \hfill
    \includegraphics[width=0.48\textwidth]{RPC_images/Velocity2.png}
    \caption{
    Gas velocity distribution in the chamber for the reference
    (left) and shifted-spacer (right) configurations.
    }
    \label{fig:rpc_velocity_comparison}
\end{figure}

The shifted-spacer configuration exhibits a more homogeneous gas
flow, reflected in a reduction of the relative velocity dispersion
from 20.3\% to 17.5\%.

\begin{figure}[htbp]
    \centering
    \includegraphics[width=0.48\textwidth]{RPC_images/Vorticity1.png}
    \hfill
    \includegraphics[width=0.48\textwidth]{RPC_images/Vorticity2.png}
    \caption{
    Vorticity distribution in the gas gap for the reference
    (left) and shifted-spacer (right) configurations.
    }
    \label{fig:rpc_vorticity_comparison}
\end{figure}

A slight reduction in bulk vorticity is observed for the optimised
layout, indicating smoother gas circulation.

\begin{figure}[htbp]
    \centering
    \includegraphics[width=0.48\textwidth]{RPC_images/Deformation1.png}
    \hfill
    \includegraphics[width=0.48\textwidth]{RPC_images/Deformation2.png}
    \caption{
    Gas-gap deformation for the reference (left) and shifted-spacer
    (right) configurations.
    }
    \label{fig:rpc_deformation_comparison}
\end{figure}

The deformation remains at the sub-percent level for both
configurations, with a small increase for the shifted layout that
remains within acceptable limits.

\begin{figure}[htbp]
    \centering
    \includegraphics[width=0.48\textwidth]{RPC_images/Electricfield1.png}
    \hfill
    \includegraphics[width=0.48\textwidth]{RPC_images/Electricfield2.png}
    \caption{
    Electric-field distribution in the gas gap for the reference
    (left) and shifted-spacer (right) configurations.
    }
    \label{fig:rpc_efield_comparison}
\end{figure}

The electric field is found to remain stable across the chamber in
both configurations, demonstrating that the optimised spacer layout
does not introduce significant field distortions.

\section{Prototype Construction and Performance Evaluation}

Two full-scale $1\times1~\si{m^2}$ glass RPC prototypes were
constructed, corresponding to the reference and shifted-spacer
configurations described in the main text. The assembly procedure
included the precise placement of spacers to define the gas gap,
graphite coating of the glass electrodes to ensure uniform surface
resistivity, sealing of the chamber with gas-tight materials, and
installation of high-voltage connections and gas inlets/outlets.

\begin{figure}[htbp]
    \centering
    \includegraphics[width=0.85\textwidth]{RPC_images/step2-1.jpg}
    \caption{Key stages of the RPC assembly.}
    \label{fig:rpc_appendix_assembly}
\end{figure}

The detector performance was evaluated using a cosmic-ray test
bench based on scintillator coincidence triggers. Two scintillator
paddles placed above and below the RPC provide a trigger signal for
through-going muons. The RPC signals are read out using front-end
electronics and digitised with a waveform digitiser for offline
analysis.

\begin{figure}[htbp]
    \centering
    \includegraphics[width=0.6\textwidth]{RPC_images/Setup.jpg}
    \caption{Cosmic-ray test setup.}
    \label{fig:rpc_appendix_setup}
\end{figure}

The detection efficiency is measured as a function of the applied
high voltage. Signal selection is performed by applying threshold
requirements on the digitised waveforms, and the stability of the
efficiency plateau is evaluated for different threshold settings.
The results are shown in Figure~\ref{fig:rpc_appendix_efficiency}.

\begin{figure}[htbp]
    \centering
    \includegraphics[width=\textwidth]{RPC_images/Sigmas.png}
    \caption{Detection efficiency as a function of high voltage for
    different signal thresholds.}
    \label{fig:rpc_appendix_efficiency}
\end{figure}

Both configurations exhibit a clear efficiency plateau above
95\%, demonstrating stable detector operation. Typical signal
waveforms and noise distributions were also analysed to characterise
the detector response and validate the choice of signal thresholds.